\subsection{Recovery in general}
\label{sec:ha:disabled:recovery}

Two requirements hold for every component in this section, and are stated once
here rather than repeated in each. Both concern the window between a component
starting and its being ready, which is short, is entered on every restart, and
is where a second failure does the most damage.

Every requirement here assumes that what survived a failure can be read. A log
truncated part-way through a record, a checkpoint that fails its checksum, or a
file that is simply absent leaves a component with nothing to recover from --- and
in this configuration with nowhere else to turn, because there is no second
instance holding a clean copy.

\textbf{This configuration is weak there, and deliberately so.} With high
availability enabled the instance whose copy is intact takes over, which is the
answer to damaged state and the reason the enabled configuration exists. Nothing
in the disabled configuration substitutes for it. That is a reason to run this
configuration only where losing what it holds is acceptable.

\begin{gap}{}
  Damaged durable state is not detected as such. A component reads what it can
  and proceeds, rather than establishing that what it read was whole, so a
  truncated log or a partial checkpoint produces a component that holds the wrong
  orders and reports itself as current. A member is then answered from that
  state: an order it placed may be unknown, and one it cancelled may be live.
\end{gap}

\begin{requirement}{R-0121}{A component takes the cost of first touching its memory before it reports itself ready}
  Where a component maps memory it will use while trading, it shall touch all of
  that memory before reporting itself ready, so that the delay of first use falls
  before it is trading rather than during.
  \rationale{Mapping a file reserves addresses and reads nothing. The first time
  each page is used the processor interrupts and the kernel fetches it --- a small
  delay, once per page, and for a region of a few hundred megabytes there are tens
  of thousands of them. Left alone they fall on whatever code touches each page
  first, which is the order path.

  Doing it before reporting ready means a component that has said it can trade
  has already paid them. Doing it afterwards moves the cost onto the first
  members to send anything, which is the worst moment to take it and the hardest
  to attribute afterwards.}
  \uncovered{no scenario measures a component's first orders against its later ones}
\end{requirement}

\begin{requirement}{R-0122}{A report a member receives twice is marked as a repeat}
  Where a recovery repeats work whose execution reports had already been sent,
  the reports sent again shall be marked so that the member can tell a repeat
  from a second event.
  \rationale{A component that dies partway through handling an order cannot know
  whether the report for it was sent. Sending it again is the right choice: a
  member that never receives one cannot recover from that alone, and a member
  that receives one twice can, provided it can tell.

  Unmarked, a repeated acceptance is indistinguishable from a second order
  accepted --- so the venue would be reporting something that did not happen, in
  the course of correcting something that did.}
  \uncovered{no scenario kills a component between sending a report and recording that it did}
\end{requirement}

\begin{requirement}{R-0031}{A component does not act on partially recovered state}
  A component shall not begin work until its recovery is complete. No sequence
  number shall be assigned before the last one issued is known, and no order
  shall be accepted before the set of open orders is current.
  \rationale{Acting on partial state produces results that look ordinary and are
  wrong. Of the ways a venue can fail, that is the hardest to detect and the
  hardest to unwind afterwards.}
  \uncovered{no scenario interrupts a recovery}
\end{requirement}

\begin{requirement}{R-0032}{An interrupted recovery can be repeated without applying anything twice}
  Where a component fails while recovering, the recovery shall be able to run
  again from the same starting point and reach the same result.
  \rationale{Without this a failure during recovery corrupts rather than delays,
  and a component that fails repeatedly leaves worse state on each attempt. It
  is also what makes a bound on the number of attempts meaningful: the bound is
  only safe if the attempts beneath it were harmless.}
  \uncovered{no scenario interrupts a recovery}
\end{requirement}
